% !TeX root = ../kazymyr-thesis.tex

\chapter{Source Code and Experiment Artifacts}\label{chapter:appendix-code}

The simulator code, the scripts that produce the charts and tables in
Chapter~\ref{chapter:experiments}, and the artifacts of all experiments are
published in the following repository.

\begin{center}
  \url{https://github.com/zenon-z/ACID-BASE-Mykhailo}
\end{center}

\noindent
The repository is private and is owned by the advisor of this thesis. Access
can be requested from the advisor.

\chapter{Use of AI Tools}\label{chapter:appendix-ai}

In line with the TUM AI Strategy, I disclose the use of generative AI tools in
this thesis. I used Claude by Anthropic, both as a chat assistant and through
Claude Code, for the following tasks.

\begin{bullets}
  \item[Learning the codebase.] I received the simulator from earlier work, and
    the tool helped me understand its structure and code faster.
  \item[Literature search.] The tool helped me find related work. I read all
    cited sources myself and checked that their content matches the way they
    are described in this thesis.
  \item[Chart scripts.] The tool helped me write the scripts that plot the
    experiment results. I reviewed the scripts and checked the charts against
    the experiment data. The scripts are available in the repository listed in
    Appendix~\ref{chapter:appendix-code}.
  \item[LaTeX tables.] The tool helped me move the results from the experiment
    sheets into LaTeX tables. I checked all values against the original data.
  \item[Proofreading.] The tool helped me correct grammar and spelling and
    improve the readability of the text.
\end{bullets}

\noindent
I reviewed all output produced with the help of AI tools, and I take full
responsibility for the content of this thesis.
